% Example Summary Report
%
% This example shows how the dwarf galaxy photometric redshift analysis should be
% written as a working scientific document.
%
% AASTeX Format Template
\documentclass[linenumbers,twocolumn]{aastex7}

\usepackage{tcolorbox}
\newtcolorbox{assumption}[1]{
  colback=blue!5!white,
  colframe=blue!75!black,
  fonttitle=\bfseries,
  title={Assumption: #1}
}

\newcommand{\sigmaNMAD}{\ensuremath{\sigma_{\rm NMAD}}}
\newcommand{\lowz}{\ensuremath{z < 0.03}}

\begin{document}

\title{Photometric Redshifts for Low-Redshift Dwarf Galaxies from DESI DR1: 
A Multi-Task Learning Approach}

\correspondingauthor{[Author Name]}

\author{First Author}
\affiliation{Stanford University}
\email{author@stanford.edu}

% Additional authors...

\begin{abstract}
We present photometric redshift estimates for low-redshift dwarf galaxies 
(\lowz) using 249,999 galaxies from DESI Data Release 1. Standard regression 
models fail catastrophically on this population due to severe class imbalance 
(only 6\% of training samples have \lowz). We address this using a multi-task 
convolutional neural network that simultaneously predicts continuous redshift 
and performs auxiliary binary/multi-class classification.

Our model achieves \sigmaNMAD\ = 0.0265 with 12.77\% outlier rate overall. 
For the target low-redshift population, outlier rates decrease from 37\% 
(baseline CNN) to 4.95\%, an 87\% reduction. We demonstrate that (1) 
end-to-end learning on images outperforms handcrafted features by a factor 
of 4.6 in scatter, (2) multi-task learning provides effective inductive bias 
for imbalanced regression, and (3) standard hierarchical CNNs appropriately 
capture multi-scale morphological information without explicit scale decoupling.

The trained model processes 9,000 galaxies per second, enabling application to 
transient host identification, dwarf galaxy surveys for dark matter studies, 
and systematic investigations of nearby galaxy evolution.
\end{abstract}

\keywords{\uat{Photometric redshifts}{1243} --- \uat{Dwarf galaxies}{416} --- 
\uat{Neural networks}{1933}}

%% ============================================================================
\section{Introduction}
%% ============================================================================

Identifying low-redshift dwarf galaxies photometrically is challenging but 
scientifically valuable. Nearby dwarfs ($z < 0.03$, corresponding to distances 
$\lesssim 130$ Mpc) provide constraints on dark matter through their dynamical 
properties and abundance \citep{Bullock2017}. They also enable time-critical 
follow-up of transients when spectroscopic host distances are unavailable 
\citep{Foley2013}.

The DESI Legacy Surveys provide deep imaging in $grz$ bands covering much of 
the sky \citep{Dey2019}. For resolved nearby galaxies, images contain distance 
information beyond integrated photometry: angular size correlates with redshift 
through the cosmological angular diameter--distance relation. Convolutional 
neural networks (CNNs) can exploit this morphological information, achieving 
substantially lower scatter than photometry-only methods 
\citep{Pasquet2019,Treyer2023}.

However, CNNs trained on the full redshift distribution systematically fail 
on low-redshift sources. The DESI LOW-Z survey \citep{DarraghFord2023} 
addresses this through sample selection and specialized classifiers, achieving 
$>$95\% completeness for $z < 0.03$. Here we take a complementary approach: 
using multi-task learning to improve regression performance on minority classes 
without explicit sample weighting.

The remainder of this paper is organized as follows. Section~\ref{sec:data} 
describes our sample selection and data characteristics. Section~\ref{sec:methods} 
presents our multi-task CNN architecture. Section~\ref{sec:results} reports 
performance metrics with comparison to baselines. Section~\ref{sec:discussion} 
discusses implications and limitations. Section~\ref{sec:conclusions} summarizes 
our key findings.

%% ============================================================================
\section{Data}
\label{sec:data}
%% ============================================================================

\subsection{Sample Selection}

We select 249,999 galaxies from DESI Data Release 1 with spectroscopic redshifts 
and $r < 21.2$. The sample combines three populations: dwarf galaxies 
($M_\star < 10^{9.25} M_\odot$), sources satisfying the DESI LOW-Z selection 
\citep{DarraghFord2023}, and a random comparison sample spanning $z < 1.5$.

\begin{assumption}{Magnitude Limit}
\textbf{Choice}: $r < 21.2$

\textbf{Rationale}: Balances sample size with image quality. At fainter 
magnitudes, photometric uncertainties increase significantly and morphological 
features become difficult to measure.

\textbf{Sensitivity}: Restricting to $r < 20.5$ reduces the sample by 40\% 
but improves \sigmaNMAD\ by $\sim$8\%.
\end{assumption}

Table~\ref{tab:sample} summarizes the redshift distribution. The target 
population ($z < 0.03$) comprises only 6\% of the sample, creating severe 
class imbalance.

\begin{deluxetable}{lcc}
\tablecaption{Redshift Distribution\label{tab:sample}}
\tablehead{
\colhead{Redshift Range} & \colhead{Count} & \colhead{Fraction}
}
\startdata
$z < 0.03$ (target) & 15,102 & 6.0\% \\
$0.03 \leq z < 0.05$ & 23,965 & 9.6\% \\
$0.05 \leq z < 0.15$ & 103,174 & 41.3\% \\
$z \geq 0.15$ & 107,673 & 43.1\% \\
\enddata
\end{deluxetable}

\subsection{Image Data}

For each galaxy, we obtain $3 \times 128 \times 128$ pixel cutouts in $grz$ 
bands from the Legacy Surveys DR9 coadded images, with pixel scale 
0.262\arcsec/pixel (33.5\arcsec\ field of view). Quality assessment on 10,000 
random images revealed 0\% NaN values, 0\% infinite values, and 0.01\% with 
extreme pixel values.

Figure~\ref{fig:sample_images} shows example low-redshift galaxies, illustrating 
the resolved morphological features---spiral arms, disk structure, irregular 
shapes---that encode distance information.

\begin{figure}
\plotone{figures/fig01_sample_images.pdf}
\caption{Example galaxies at $z < 0.03$ showing $grz$ bands and RGB composite. 
Resolved internal structure distinguishes nearby galaxies from more distant 
compact sources.\label{fig:sample_images}}
\end{figure}

%% ============================================================================
\section{Methods}
\label{sec:methods}
%% ============================================================================

\subsection{Multi-Task Architecture}

We use a ResNet-18 backbone \citep{He2016} with three prediction heads:
\begin{enumerate}
    \item \textbf{Regression}: Continuous redshift via fully-connected layer
    \item \textbf{Binary classification}: $z < 0.03$ vs.\ $z \geq 0.03$
    \item \textbf{Multi-class classification}: Low/mid/high-$z$ bins
\end{enumerate}

The combined loss is
\begin{equation}
\mathcal{L} = \mathcal{L}_{\rm Huber} + \lambda_1 \mathcal{L}_{\rm BCE} + 
\lambda_2 \mathcal{L}_{\rm CE}
\end{equation}
where $\mathcal{L}_{\rm Huber}$ is the Huber loss for regression ($\delta = 1.0$), 
$\mathcal{L}_{\rm BCE}$ is binary cross-entropy, $\mathcal{L}_{\rm CE}$ is 
categorical cross-entropy, and $\lambda_1 = \lambda_2 = 0.5$.

\begin{assumption}{Multi-Task Weighting}
\textbf{Choice}: $\lambda_1 = \lambda_2 = 0.5$

\textbf{Rationale}: Equal weighting provides balanced contribution from auxiliary 
tasks without dominating the regression objective.

\textbf{Sensitivity}: $\lambda = 0.1$ (weaker auxiliary signal) increases low-$z$ 
outliers by 5\%. $\lambda = 1.0$ (stronger signal) slightly degrades overall 
\sigmaNMAD.
\end{assumption}

The auxiliary classification tasks provide inductive bias that encourages the 
network to learn features discriminative for the minority low-$z$ class. Without 
explicit class weighting, the regression loss is dominated by the majority 
populations.

\subsection{Training}

We split the data 70/15/15 into training (174,999), validation (37,500), and 
test (37,500) sets using stratified sampling to preserve class proportions. 
Training uses Adam optimizer with learning rate $10^{-3}$, batch size 256, and 
ReduceLROnPlateau scheduler (patience 5, factor 0.5). Training runs for 50 
epochs with early stopping based on validation \sigmaNMAD.

Data augmentation includes random horizontal/vertical flips and 90$^\circ$ 
rotations, which are valid because sky orientation carries no distance information. 
Per-image normalization ($x_{\rm norm} = (x - \mu)/\sigma$) accounts for varying 
observing conditions.

%% ============================================================================
\section{Results}
\label{sec:results}
%% ============================================================================

\subsection{Overall Performance}

Table~\ref{tab:results} compares model performance on the held-out test set. 
Our Multi-Task CNN achieves \sigmaNMAD\ = 0.0265 with 12.77\% outlier rate 
overall, and critically reduces low-$z$ outliers from 37\% (baseline) to 4.95\%.

\begin{deluxetable}{lcccc}
\tablecaption{Model Performance on Test Set\label{tab:results}}
\tablehead{
\colhead{Model} & \colhead{\sigmaNMAD} & \colhead{Outliers} & 
\colhead{Low-$z$ Outliers} & \colhead{Bias}
}
\startdata
XGBoost (features) & 0.042 & 27.0\% & 31.1\% & +0.003 \\
Baseline CNN & 0.036 & 18.2\% & 37.4\% & +0.012 \\
ResNet-18 & 0.028 & 14.6\% & 11.0\% & +0.008 \\
\textbf{Multi-Task CNN} & \textbf{0.027} & \textbf{12.8\%} & \textbf{5.0\%} & +0.005 \\
\enddata
\tablecomments{Outliers defined as $|\Delta z / (1+z_{\rm true})| > 0.05$. 
Best performance in bold.}
\end{deluxetable}

\begin{figure}
\plotone{figures/fig02_predictions.pdf}
\caption{Predicted vs.\ true redshift for the Multi-Task CNN (N = 37,500). 
Gray points show all galaxies; red highlights the $z < 0.03$ population. 
The dashed line indicates perfect agreement. The model achieves \sigmaNMAD\ = 
0.0265 overall and 0.0123 for $z < 0.03$.\label{fig:predictions}}
\end{figure}

\subsection{Low-Redshift Performance}

For the target $z < 0.03$ population, the Multi-Task CNN achieves:
\begin{itemize}
    \item \sigmaNMAD\ = 0.0123
    \item Outlier rate: 4.95\% (vs.\ 37.4\% baseline)
    \item Binary classification ROC-AUC: 0.963
\end{itemize}

The 87\% reduction in low-$z$ outliers relative to baseline CNNs represents 
the primary contribution of this work.

\subsection{Comparison to Feature Engineering}

Direct comparison reveals that end-to-end learning substantially outperforms 
handcrafted features: \sigmaNMAD\ improves by a factor of 4.6 (0.122 $\to$ 0.027) 
and outlier rates by 5.3$\times$ (67\% $\to$ 13\%). Feature-based models showed 
severe overfitting (validation \sigmaNMAD\ $\approx$ 0.03 vs.\ test \sigmaNMAD\ = 
0.12), while CNNs generalized well (validation $\approx$ test performance).

%% ============================================================================
\section{Discussion}
\label{sec:discussion}
%% ============================================================================

\subsection{Why Multi-Task Learning Works}

The auxiliary classification heads force the network to learn features 
discriminative for the minority low-$z$ class. Without explicit class weighting, 
the regression loss is dominated by the majority mid-$z$ and high-$z$ populations. 
The classification heads provide an inductive bias that prioritizes separating 
the low-$z$ population, and these features transfer to improve regression 
performance.

This insight generalizes: for imbalanced regression problems, auxiliary 
classification tasks on binned targets provide an alternative to sample 
weighting or loss function modifications.

\subsection{Limitations}

Several limitations should be noted:

\begin{enumerate}
    \item \textbf{Single survey}: All results are from DESI DR1. Performance on 
    other surveys (SDSS, HSC, Euclid, Rubin LSST) is unknown and may differ due 
    to varying depth, resolution, and filter sets.
    
    \item \textbf{Three bands only}: We use only $grz$ photometry. Additional 
    bands (especially infrared from WISE) could improve high-$z$ performance 
    where optical colors provide limited information.
    
    \item \textbf{Point predictions}: We provide no uncertainty estimates. The 
    classification probabilities could serve as proxies, but proper uncertainty 
    quantification requires additional methodology (e.g., ensembles, probabilistic 
    outputs).
    
    \item \textbf{Systematic biases}: All models show redshift-dependent biases: 
    overprediction at low-$z$ (bias = +0.014) and underprediction at high-$z$ 
    (bias = $-$0.011). Post-hoc calibration may be needed for applications 
    sensitive to these biases.
\end{enumerate}

\subsection{Implications for Transient Astronomy}

The model processes 9,000 galaxies per second on GPU, enabling real-time host 
galaxy classification for transient alerts. For a candidate supernova with no 
spectroscopic host redshift, applying this model to the Legacy Survey cutout 
can provide a distance estimate within seconds.

A binary classification threshold of 0.3 achieves 70\% recall for true low-$z$ 
hosts with 45\% purity---sufficient for prioritizing spectroscopic follow-up.

%% ============================================================================
\section{Conclusions}
\label{sec:conclusions}
%% ============================================================================

We presented a Multi-Task CNN for photometric redshift estimation of low-redshift 
dwarf galaxies from DESI DR1. The key findings of this work are:

\begin{enumerate}
    \item Multi-task learning with auxiliary classification heads reduces 
    low-$z$ outliers by 87\% compared to baseline CNNs (from 37\% to 5\%).
    
    \item End-to-end deep learning outperforms handcrafted feature engineering 
    by a factor of 4.6 in scatter.
    
    \item Standard hierarchical CNNs appropriately capture multi-scale 
    morphological features; explicit scale-morphology decoupling degrades 
    performance.
    
    \item The model achieves \sigmaNMAD\ = 0.0265 overall with inference speed 
    of 9,000 galaxies/second, suitable for billion-galaxy surveys.
\end{enumerate}

Future work should explore uncertainty quantification through ensemble methods, 
transfer learning to other imaging surveys, and extension to additional 
photometric bands.

\begin{acknowledgments}
We thank [collaborators] for helpful discussions. This work used data from 
DESI Data Release 1. [Standard DESI acknowledgment]
\end{acknowledgments}

\facilities{DESI, Mayall (DESI)}

\software{
astropy \citep{astropy2022},
matplotlib \citep{matplotlib},
numpy \citep{numpy},
pytorch \citep{pytorch},
scikit-learn \citep{sklearn}
}

\bibliography{ms}

\end{document}
